\section{Fundamentos de las ecuaciones diferenciales estocásticas (SDE)}

\subsection{Forma general de las SDE}

En el dominio de tiempo continuo, el proceso de difusión hacia adelante de los modelos de difusión puede representarse mediante la siguiente SDE:

$$
\mathrm{d}\mathbf{x} = f(\mathbf{x}, t)\,\mathrm{d}t + g(t)\,\mathrm{d}\mathbf{w}
$$

Donde cada símbolo tiene el siguiente significado:

\begin{itemize}
    \item $\mathbf{x}$ —— Variable de datos (por ejemplo, imágenes), que es un proceso estocástico en función del tiempo $t$;
    \item $f(\mathbf{x}, t)$ —— \textbf{coeficiente de deriva} (drift coefficient), que determina la dirección de evolución temporal de la media de $\mathbf{x}_t$;
    \item $g(t)$ —— \textbf{coeficiente de difusión} (diffusion coefficient), que controla la intensidad del ruido inyectado;
    \item $\mathbf{w}$ —— \textbf{proceso de Wiener estándar} (proceso de Wiener), también conocido como movimiento browniano;
    \item $\mathrm{d}\mathbf{w}$ —- incremento infinitesimal del proceso de Wiener.
\end{itemize}

\begin{importantbox}{Intuición sobre la deriva y la difusión}
El \textbf{término de deriva} $f(\mathbf{x}, t)\,\mathrm{d}t$ puede entenderse como una "fuerza" determinista que mueve cada punto de datos en cierta dirección (por ejemplo, contrayéndose hacia el origen). El \textbf{término de difusión} $g(t)\,\mathrm{d}\mathbf{w}$ inyecta ruido aleatorio, haciendo que la trayectoria de $\mathbf{x}_t$ se vuelva aserrada. La combinación de ambos produce un mapeo continuo desde la distribución de datos hasta la distribución de ruido.
\end{importantbox}

\subsection{Proceso de Wiener (movimiento browniano)}

El proceso de Wiener $\mathbf{w}(t)$ es la fuente de aleatoriedad en las SDE y posee las siguientes cuatro propiedades características:

\begin{enumerate}
    \item \textbf{Condición inicial}: $\mathbf{w}(0) = 0$ (casi seguro, es decir, $P(\mathbf{w}(0)=0) = 1$);
    \item \textbf{Normalidad de los incrementos}: para cualquier $0 \le s < t$, el incremento $\mathbf{w}(t) - \mathbf{w}(s) \sim \mathcal{N}(0, (t-s)\mathbf{I})$;
    \item \textbf{Incrementos independientes}: los incrementos en intervalos de tiempo no superpuestos son mutuamente independientes;
    \item \textbf{Continuidad de las trayectorias}: $\mathbf{w}(t)$ es continuo respecto a $t$ casi seguro.
\end{enumerate}

\begin{knowledgebox}{El propio proceso de Wiener es una SDE}
Si se establece $f = 0$ y $g = 1$, la SDE se reduce a $\mathrm{d}\mathbf{x} = \mathrm{d}\mathbf{w}$, cuya solución es el propio proceso de Wiener. Por lo tanto, el movimiento browniano puede considerarse la SDE más simple: no hay deriva, solo difusión puramente aleatoria.
\end{knowledgebox}

\subsection{Desarrollo en serie de Taylor y aproximación}

Al convertir el proceso de difusión discreto a una SDE continua, el desarrollo en serie de Taylor es la herramienta matemática fundamental. Dada una función suave $f$, su aproximación de primer orden es:

$$
f(t + \Delta t) \approx f(t) + f'(t) \cdot \Delta t
$$

Por ejemplo, para $f(\epsilon) = \sqrt{1 - \epsilon}$, cuando $\epsilon$ es muy pequeño:

$$
\sqrt{1 - \epsilon} \approx 1 - \frac{\epsilon}{2}
$$

Esta aproximación se utiliza repetidamente al deducir las funciones $f(t)$ y $g(t)$ del VP-SDE. Cuando el tamaño del paso temporal $\Delta t \to 0$, los términos de orden superior desaparecen, obteniendo así la forma exacta de la ecuación diferencial.

\begin{warningbox}{Supuestos previos a la aproximación continua}
La conversión de lo discreto a lo continuo depende de las hipótesis de que el tamaño del paso temporal es suficientemente pequeño ($\Delta t \to 0$) y que la función de programación del ruido (como $\beta_t$) es suficientemente suave. Para la discretización con un número finito de pasos en implementaciones prácticas, la SDE continua es solo una aproximación: a mayor número de pasos, la aproximación será más precisa. Esto también explica por qué aumentar el número de pasos de discretización suele mejorar la calidad generada.
\end{warningbox}

\subsection{Resumen de la sección}

Las SDE describen la evolución aleatoria de los datos en tiempo continuo a través de los términos de deriva y difusión. El proceso de Wiener proporciona el impulso estocástico estándar. Mediante el desarrollo en serie de Taylor, podemos partir de las probabilidades de transición de una cadena de Markov discreta y deducir la SDE continua correspondiente.
